\section{Introduction}
\label{sec:intro}

For graph transformation tools to be able to model practically relevant domains, they have to support a notion of data attributes. In the case of \GROOVE, this support is rooted in algebraic theory: there is a pre-defined signature with four primitive sorts (\intS, \realS, \boolS and \stringS), with for each sort a set of constants and a collection of basic operators. \GROOVE offers a choice of four different algebras at various levels of abstraction (the \termA algebra, in which no computation takes place; the \bigA algebra for infinite-precision arithmetic; the \javaA algebra, directly based on Java types with \doubleT for \realS and \intT for \intS; and the \pointA algebra, where each carrier sort contains precisely one value).
%
Data values are incorporated in graphs as special \emph{value nodes}, typed by their sorts; attributes are edges from plain (non-value) nodes to value nodes. Rules can contain sort-typed variable nodes and expressions built from operators and variable nodes.

This setup has clear limitations. For instance, it is awkward to model advanced data manipulation beyond those pre-defined types and operations, such as trigonometrical or statistical calculations. Moreover, the algebraic theoretical underpinning means that operations are \emph{determinate}: their outcome is always completely determined by their arguments. This means that, for instance, randomness cannot be modelled. Finally, there is no support for structured (value) data types, such as records; while this does not affect the modelling power per se, it does hamper the usability for practical purposes.

\medskip\noindent
In this paper, we describe an extension of \GROOVE to \emph{used-defined types and operations}, based on programmable Java classes and methods, that gets rid of these limitations. Essentially, the user can provide annotated record classes and classes with annotated static methods; by making these available to the \GROOVE library at link time, they can be used in rule systems. User-defined record types are limited to primitively typed fields; user-defined operations can have any signature based on the four primitive types plus user types. Moreover, methods can be declared as \emph{indeterminate}, meaning that their outcome is not completely determined by their arguments. This allows the inclusion of randomness and other forms of non-determinism in data manipulation.

Below, we briefly go over the theoretical underpinning of \GROOVE attributes (\Cref{sec:background}) and show how the user-defined types and operations fit into that; then we show the actual implementation on the basis of an example (\Cref{sec:example}). Related work and conclusions are presented in \Cref{sec:concl}.

The functionality described in this paper has been incorporated in \href{https://github.com/nl-utwente-groove/code/releases/tag/release-\GROOVEversion}{\GROOVE release~\GROOVEversion}, and the Java code and \GROOVE rule system of the example are available as part of \href{https://github.com/nl-utwente-groove/userops.git}{this paper's repository}.
